%\DeclareGraphicsExtensions{.png}
\chapter{Redox Reactions}
\def\shorttitle{Redox Reactions} 
\section[Background]{Background}

Understanding oxidation-reduction (redox) reactions that occur in natural 
and contaminated systems is central to developing accurate conceptual and 
quantitative models. Redox reactions involve the transfer of electrons 
between constituents and therefore, changes in the oxidation state of 
one or more elements. Changes in oxidation state of an element can 
result in dramatic changes in its chemical and physical properties.  
For example, iron in the plus three oxidation 
state - Fe(III) - in surficial deposits renders 
a reddish hew to Mars whereas iron in the plus II oxidation 
state - Fe(II) - in surficial basalts on the Moon is, in no 
small part, responsible for the blackness of the lunar mares. 
Two important examples of redox reactions are the oxidation 
of aqueous Fe(II) to Fe(III) or organic carbon by dissolved oxygen:

\begin{equation}
4Fe^{+2} + O_{2,aq} + 10H_2O \rightarrow 4Fe(OH)_{3,s} + 8H^+  \\
\label{eq:fe2_oxidation}
\end{equation}

\begin{equation}
CH_2O  + O_{2,aq} + 10H_2O \rightarrow CO_{2,aq} + H_2O \\
\label{eq:om_reduction}
\end{equation}

Rather than representing formaldehyde, $CH_2O$ is often used to represent 
one sixth of a glucose molecule.

The classical approach to discussing redox reactions in aqueous geochemistry 
is to split the redox reactions into two hypothetical 
half reactions involving electrons. 
Using this approach, reaction \ref{eq:fe2_oxidation} and \ref{eq:om_reduction} 
would be broken up into the following half reactions:

\begin{equation}
Fe^{+2} + 3H_2O  \rightarrow Fe(OH)_{3,s} + 3H^+ + e^-  \\
\label{eq:fe_half_reac}
\end{equation}

\begin{equation}
CH_2O + H_2O \rightarrow CO_{2,aq} + 4H^+ + 4e^-  \\                    
\label{eq:om_half_reac}
\end{equation}

\begin{equation}
O_{2,aq} + 4H^+ + 4e^-  \rightarrow 2H_2O \\
\label{eq:ox_half_reac}
\end{equation}

Mass action expressions are then written for the half reactions.
For example, for the Fe(II)/Fe(III) half reaction, 
the mass action expression is:

\begin{equation}
K = a_{Fe(OH)_{3,s}} a_{H^+}^3 a_{e^-}/a_{Fe}^{+2} a_{H_2O}^3 \\
\label{eq:rdx6}
\end{equation}

The standard state of the $e^-$ is the same as other aqueous species,
1 mole per kilogram water.  
This approach comes from the field of electrochemistry, where the oxidation and 
reduction half reactions can be made to occur in separate place
connected, for example, 
through a meter that measures the electromotive force (or power source) 
and a salt bridge.
In dilute solutions the activity of water can be equated with its mole fraction,
which is essentially 1. If pure ferric hydroxide forms then it 
can be assumed to be in 
its standard state (mole fraction $Fe(OH)_{3,s}$
equals 1), 
so Eqn (\ref{eq:rdx6}) can be re-written:

\begin{equation}
K=a_{H^+}^3 a_{e^-}/a_{{Fe}^{+2}} \\    
\label{eq:rdx7}
\end{equation}
                                                         
Taking the base-10 logarithm, \ref{eq:rdx7} becomes: 
\begin{equation}
logK = 3loga_{H^+} + loga_{e^-} -loga_{{Fe}^{+2}}  \\                          
\label{eq:rdx8}
\end{equation}

Analogous to the definition of $pH$,
$loga_{e^-}$
is defined as $pe$.
This formalism is useful for graphically representing the 
relative thermodynamic stability of redox-reactive species 
in natural and contaminated systems 
(typically in their standard state).

Electrochemists have established logK values for 
half reactions based on the 
convention that the hydrogen gas -
hydrogen ion half reaction
\begin{equation}
H_{2,g}  =  2H^+ + 2e^- \\   
\label{eq:rdx9}
\end{equation} 
has an electromotive force and, therefore a logK 
value of 0 for $H_{2,g}$ and $H^+$ in their 
standard states (1 atmosphere and 1 mol/kgw, respectively).  
Using this convention, logK values for half reactions can be assembled 
in order of thermodynamic favorability, 
called a redox ladder, such as Table \ref{tab:rdx1}.  
Note that the logK values apply to all species in their standard state 
(1 mol/kgw for aqueous species, pure solid phase activity 
equals 1, and 1 atmosphere for gasses).  
Also, the logK value for half reactions where solid phases are participants 
depends on the thermodynamic stability of the solid phase.  
For example, reduction of $ferrihydrite$ (fhy) is more favorable than 
reduction of $goethite$ (based on the logK values for goethite that 
applies to the Cape Cod aquifer, Kent et al., 2008) because $goethite$ 
is thermodynamically more stable than $ferrihydrite$.  
Plots of $pH$ against $pe$ constructed using reactions and logK values like 
those in Table \ref{tab:rdx1} are used to illustrate 
trends expected in redox processes 
with changes in pH. In the exercises that will follow below, 
we will explore these trends.

\begin{table}[h]
\topcaption{Half reactions and logK values for some redox-reactive species}
\centering
%\footnotesize
\begin{tabular}{rcll}
\toprule
Half reactions, per electron				  & &									                  & logK \\ \hline
0.25 $O_{2,aq}$ + $e^-$ + $H^+$       &=& $0.5H_2O$											    & 21.52 \\
0.2  $NO_3^-$ + $e^-$ + $1.2H^+$      &=& 0.1 $N_{2,aq}$ + 0.6 $H_2O $      & 20.71 \\
0.5 $MnO_2:H_2Os$ + $e-$ + $2H^+$     &=& $0.5Mn^{+2}$ + $1.5H_2O$				  & 20.69 \\
0.5 $HAsO_4^{-2}$ + $e^-$ + $2H^+$    &=&  0.5 $As(OH)_{3,aq}$ +  $0.5H_2O$	      & 19.80 \\
$FeOOH$ + $e^-$ + $3H^+$              &=& $Fe^{+2}$ + $2H_2O$  fhy							    & 17.91 \\
$FeOOH$ + $e^-$ + $3H^+$              &=& $Fe^{+2}$ + $2H_2O$  Goethite (Kent et al. 2008)	    & 13.52 \\
$0.125SO_4^{-2}$ + $e^-$ + $1.125H^+$ &=& $0.125HS^-$ + $0.5H_2O$							    & 4.21 \\
$0.125CO_3^{-2}$ + $1.25H^+$ + $e^-$  &=& $0.125CH_{4,aq}$ + $0.375H_2O$						& 5.13 \\
\bottomrule
\label{tab:rdx1}
\end{tabular}
\end{table}	

\begin{table}[h]
\topcaption{PHREEQC keywords for simulating redox equilibria}
\centering
\footnotesize
\begin{tabular}{llll}
\toprule
Constituents	& Database	                & Input	     & Selected\_Output   \\ \hline
Aqueous       &	SOLUTION\_MASTER\_SPECIES & SOLUTION n & Totals (e.g., Ca)   \\
              & SOLUTION\_SPECIES	        & MIX n      & Molalities (e.g., $Ca^{+2}$) \\
              &                           &            & Activities (e.g., $Ca^{+2}$) \\
              &                           &            & Saturation\_index  \\
              &                           &            & USER\_PUNCH       \\
              &                           &            &        SC         \\
              &                           &            &     Rho           \\
              &                           &            &   Soln\_Vol        \\ \hline
Gas phase	    & Phases	   & EQUILIBRIUM\_PHASES n   & EQUILIBRIUM\_PHASES \\
              &            &                         & USER\_PUNCH   \\
              &            &                         & EQUI(...)     \\ \hline
Solid phase	  & Phases     & EQUILIBRIUM\_PHASES n   & See Gas Phase \\ \hline
Surface       & SURFACE\_MASTER\_SPECIES & SURFACE n & See aqueous   \\
complexation  & SURFACE\_SPECIES         &           &       \\ \hline
Other         &                          &           &       \\ \hline
SAVE, USE	    & Solution n,              &           &       \\
              & EQUILIBRIUM\_PHASES n    &           &       \\ 
							& EXCHANGE  n              &           &       \\ 
							& SURFACE n                &           &       \\ \hline
REACTION n    &	Irreversible addition    &           &       \\
              & of compounds             &           &       \\ 

\bottomrule
\label{tab:keywords}
\end{tabular}
\end{table}	

\section[PHREEQC Exercise 6]{PHREEQC Exercise 6}

This series of exercises involves using PHREEQC to react 
oxidized and reduced species. As you will see, 
the oxidation reduction reactions 
also induce changes in pH and other secondary reactions.


\begin{table}[h]
\caption{Water compositions for the mixing exercise. Adjust Na to achieve charge balance for the oxic solution and Cl to achieve charge balance for the anoxic solution.}
\begin{center}
\begin{tabular}{lccc} 
\toprule
Parameter   & Oxic GW & Anoxic GW & Units  \\ \hline
pH          & 6.50 		& 6.50 			& -      \\
pe          & 12   		& 4    			& -      		\\
Temperature & 21   		& 21   			& $^\circ$C \\
Alkalinity  & 710  		& 710  			& $\mu$eq/L  		\\
Na          & 780  		& 580  			& $\mu$mol/L \\ 
K           & 8.0  		& 75   			& $\mu$mol/L \\
Mg          & 18.0 		& 53   			& $\mu$mol/L \\
Ca          & 9.0  		& 170  			& $\mu$mol/L \\ 
Cl          & 82   		& 550  			& $\mu$mol/L \\ 
S(6)        & 24.0 		& 100  			& $\mu$mol/L \\ 
O(0)        & 500  		& 0    			& $\mu$mol/L \\ 
Fe(2)       & 0    		& 179 			& $\mu$mol/L \\ 
\bottomrule
\end{tabular}
\label{tab:rdx2}
\end{center}
\end{table}

\subsection[Mixing-induced redox reactions]{Mixing-induced redox reactions}

\begin{itemize}
\item Set up two solutions, one with $O_2$ and one with $Fe(II)$, using the compositions in Table \ref{tab:rdx2}. {\color{red} \textbf{MIX}} them in equal proportions. (Keywords: {\color{red} \textbf{SOLUTION n}}, {\color{red} \textbf{SOLUTION m}}, {\color{red} \textbf{MIX o}}). The solutions begin with the same $pH$ and $alkalinity$. What about the $TCO_2$ concentration ? $TCO_2$ is the sum of the concentrations of all carbonate-containing aqueous species. 
	\begin{enumerate} 
	\item What happens to the $pH$, $alkalinity$, and $TCO_2$ concentrations after mixing ?
	\item What happens to the concentrations of $O_2$, $Fe(II)$, and $Fe(III)$ after mixing?  What are the dominant $Fe$ aqueous species ?
	\item What happens to the $PCO_2$ after mixing?
	\end{enumerate} 
\item Use the anoxic solution. Allow 10 umoles of $O_2(g)$ from the atmosphere ($PO_2$ = 0.2 atmospheres) to equilibrate with the solution. 
(Keyword: {\color{red} \textbf{EQUILIBRIUM\_PHASES n}}.)
  \begin{enumerate} 
	\item What happens to the $pH$, $alkalinity$, and $TCO_2$ concentrations after mixing ?
	\item What happens to the concentrations of $O_2$, $Fe(II)$, and $Fe(III)$ after mixing?  What are the dominant $Fe$ aqueous species?
	\item What happens to the $PCO_2$ after mixing?
  \end{enumerate}
\item Repeat the previous exercise where first 100 $\mu$moles then 1000 $\mu$moles of $O_2$ from the atmosphere equilibrate with the anoxic solution, but allow ferrihydrite ($Fe(OH)3(a)$) or goethite ($Goethite$) to precipitate until it reaches equilibrium.
	\begin{enumerate}  
	\item What happens to the $pH$, $alkalinity$, and $TCO_2$ concentrations after mixing ?
	\item What happens to the concentrations of $O_2$, $Fe(II)$, and $Fe(III)$ after mixing?  What are the dominant $Fe$ aqueous species ?
	\item What happens to the $PCO_2$ after mixing ?
	\item What happens when the solution is allowed to equilibrate with atmospheric $PCO_2$ ?
	\end{enumerate} 
\item Simulate reacting the oxic solution with 500 umoles of amorphous $FeS$ ($FeS(ppt)$) (ca. 44 mg) allowing $Fe(OH)3(a)$ to precipitate. 
	\begin{enumerate}  
	\item What got oxidized ? What did not get oxidized? What was the $pH$ ?
	\item Repeat the simulation, this time allowing the suspension to equilibrate with atmospheric $O_2$. What was the $pH$ ? What got oxidized ?
	\end{enumerate}  
\end{itemize} 


\subsection[Redox sequence]{Redox sequence}

The next exercise compares the standard sequence of oxidation- reduction reactions 
predicted using the redox ladder to the sequence in which oxidants (electron acceptors) 
actually get reduced by organic carbon based on thermodynamic considerations.   
The exercise uses the keyword {\color{red} \textbf{REACTION n}} to add $CH_2O$, 
which represents one six$^{th}$ of a glucose molecule, to a suspension with aqueous and solid-phase oxidants.  
As $CH_2O$ is added it reacts with the oxidants that are most thermodynamically favorable. 
Thermodynamic favorability depends on concentrations and solid-phase solubilities 
(which depend on $pH$ for some important oxidants).  
Thus the order in which the oxidants react with $CH_2O$ may change during the course 
of the redox titration as the chemical conditions change.  

\textbf{Note, that you will need to use either the {\color{red} \textbf{Wateq4f.dat}} 
or {\color{red} \textbf{Minteq.v4.dat}} database, 
because the Phreeqc.dat database does not include arsenic or uranium.}


\begin{itemize}
\item Create a PHREEQC input file and set up the initial solution and solid-phase compositions 
in the suspension using the data in Table \ref{tab:rdx3}.  
Assume the suspension has one kilogram of water.  
(Keywords: {\color{red} \textbf{SOLUTION n}}, {\color{red} \textbf{EQUILIBRIUM\_PHASE n}}).  
Adjust the chloride concentration to insure charge balance of the solution. 
Two solid phases with reduced species may precipitate if they become supersaturated: 
amorphous $FeS$ ($FeS(ppt)$ in Phreeqc) and $Siderite$.

\begin{table}[h]
\caption{The composition of the suspension used for the redox sequence and redox titration exercises}
\begin{center}
\begin{tabular}{lcc} 
\hline
	   & Value	 & Units \\
pH	 & 5.80	   & - \\
T 	 & 22	     & $^\circ$C \\
C(4) &	2514	 & $\mu$M \\
Na	 & 2580	   & $\mu$M \\
K	   & 330	   & $\mu$M \\
Mg	 & 170	   & $\mu$M \\
Ca	 & 570	   & $\mu$M \\
Cl	 & 1766	   & $\mu$M \\ 
Water	& 1	     & Kg = L \\ \hline % should be approx. not equal
Redox-reactive constituents & &\\ \hline
O2 	  & 250	      & $\mu$M \\
N(5)	& 1500	    & $\mu$M (use Nit not N) \\
Pyrolusite & 800	& $\mu$moles \\
As(5)      &	1	  & $\mu$M \\
U(6)	     &  1	  & $\mu$M \\
Fe(OH)3(a) or Goethite 	& 8000	& $\mu$moles \\
S(6) 	                  & 260	  & $\mu$M \\ \hline
\end{tabular}
\label{tab:rdx3}
\end{center}
\end{table}


Using thermodynamic data from the PHREEQC database 
(and the Wateq4f or Minteq.v4 database if arsenic or uranium are included), 
the redox ladder for the oxidized species in Table \ref{tab:rdx1} follows the sequence: 
$O_2$, $NO_3^-$, $MnO_2$ ($Pyrolusite$), $As(V)$, ferric hydroxide (ferrihydrite, $Fe(OH)3(a)$ 
in the PHREEQC database, or $Goethite$), and $SO_4^{-2}$ 
(Table \ref{tab:rdx1}).  
This sequence is based on all reactants in Table \ref{tab:rdx1} 
being in their standard state at $pH$ 7 and 25 $^\circ$C.  
Use the {\color{red} \textbf{SAVE}} keyword to save 
the composition of the solution and the equilibrium phase assemblage.

\item Use the keyword {\color{red} \textbf{USE}} to recall the compositions 
of the solution and equilibrium phase assemblage created previously.  
Use the {\color{red} \textbf{REACTION n}} keyword to add an increment of $CH_2O$.  
Try adding a suffiently large amount of $CH_2O$ to remove about half of the dissolved oxygen, 
then try again with a now sufficiently large amount of $CH_2O$ 
to remove all the dissolved oxygen and even some of the nitrate, etc.  
Examine the order in which oxidants get reduced. Any surprises ?

\item Using the keyword {\color{red} \textbf{SELECTED\_OUTPUT}}, 
set up a file that includes simulation data of interest to this exercise 
without a lot of the other information that, while interesting, 
clutters the most significant simulation results.  
Consult the manual under the {\color{red} \textbf{SELECTED\_OUTPUT}} 
keyword to have the following information sent to the selected output file: 
$pH$, $alkalinity$, concentrations of the total dissolved oxygen ($O(0)$), nitrate ($N(5)$), 
$Pyrolusite$, $As(5)$, ferric hydroxide ($Fe(OH)3(a)$ or $Goethite$, 
whichever you chose to use), $S(6)$, $Siderite$, and $FeS(ppt)$.  
For the solid phases, you can use the {\color{red} \textbf{-equilibrium\_phases}} 
list-of-phases command under {\color{red} \textbf{SELECTED\_OUTPUT}}.  
For each phase in the list, this will send the total number of moles 
remaining and the number of moles added (positive value) or removed (negative value) 
to the {\color{red} \textbf{selected output file}}.  
Alternatively, you can also add a {\color{red} \textbf{USER\_PUNCH}} 
list to your {\color{red} \textbf{SELECTED\_OUTPUT}}.  
Consult the PHREEQC manual under the {\color{red} \textbf{USER\_PUNCH}} keyword.  
It will explain how to add the appropriate heading.
To add the moles of a solid (or gas) phase, use 
{\color{red} \textbf{PUNCH Phases (phase name)}}, % shoul be Phases("`phase name”)
as explained in Table 8 of the PHREEQC manual.  
The selected output file (note, you should not forget to give it a name 
using the {\color{red} \textbf{-filename}} 
name command) can be opened using the Excel wizard 
(it as a delimited file).

\end{itemize}

\subsection[Redox titration]{Redox titration}

In this exercise we conduct a redox titration, where $CH_2O$ 
is added incrementally to a suspension with aqueous- and solid-phase oxidants.  
We will begin by assuming that the oxidants will be reduced sequentially based 
on the redox ladder in Table \ref{tab:rdx1} and compare this to the actual sequence.
The exercise is broken down into a number of steps:

\begin{enumerate}
\item Determine how many moles (per liter) of $CH_2O$ to you need add to titrate away each of the oxidants. One possible approach is described below.
\item Set up the solution and phase compositions. The beginnings of a PHREEQC file is provided (in an Excel file). To that file, add the solution composition and equilibrium phase assemblage in Table Table \ref{tab:rdx1}.
\item Set up a {\color{red} \textbf{SELECTED\_OUTPUT}} file following instruction in step 3 of exercise 5, above. 
\item Using the {\color{red} \textbf{REACTION keyword}}, 
add blocks to the input file in which each oxidant is titrated away sequentially be $CH_2O$ in a series of small steps.
\item Run the simulation, plot the results.
\end{enumerate}

\begin{itemize}
\item Determine how many moles of $CH_2O$ must be added to 
titrate away each of the oxidants in Table \ref{tab:rdx1}. 
One approach is first to establish the moles of $CH_2O$ required 
to reduce one mole of each oxidant and second to multiply the concentrations 
of each oxidant by the ratio of $CH_2O$ reduced per mole of oxidant.
\end{itemize}

The half reaction for the oxidation of $CH_2O$ to $CO_2$ written:



\begin{equation}
CH_2O  + H_2O \rightarrow CO_{2,aq} + 4H^+ + 4e- \\
\label{eq:half_reac_om}
\end{equation}

Four electrons are released by the reduction of one $CH_2O$.  
The half reaction for reduction of $O_2$ (Table \ref{tab:rdx1}) is written:

\begin{equation}
O_{2,aqu}  + 4e- + 4H^+ \rightarrow 2H_2O \\
\label{eq:half_reac_ox}
\end{equation}

Four electrons are consumed for each $O_2$ reduced.  
So the ratio of $CH_2O$ oxidized to $O_2$ reduced is 1.  
In order to reduce 250 $\mu$moles (per liter) of dissolved oxygen, 
250 $\mu$moles (per liter) of $CH_2O$ must be added.
Repeating this process allows a set of balance redox 
reactions to be written, from which the stoichiometric relationship 
between $CH_2O$ added and oxidant removed can be determined 
(Table \ref{tab:rdx4}).


\begin{table}[h]
\topcaption{Overall reactions for the oxidation of $CH_2O$ by dissolved oxygen, nitrate, manganese oxide, arsenate, FeOOH, and sulfate.}
\centering
\footnotesize
\begin{tabular}{lcr}
\toprule
Overall reaction		& &														         \\ \hline 
$O_{2,aq}$ +  $CH_2O$               &=& $CO_{2,aq}$ + $H_2O$ = $HCO_3^-$ + $H^+$  \\
4 $NO_3^-$ +   $5CH_2O$ + $4H^+$    &=&  2 $N_{2,aq}$ + 5 $CO_{2,aq}$ + 7 $H_2O$ \\
2 $MnO_2·H_2Os$ + $CH_2O$ + 4 $H^+$ &=& 2 $Mn^{+2}$ + $CO_{2,aq}$ + 5 $H_2O$ \\
2 $H_2AsO^{4-}$ + $CH_2O$ + 2 $H^+$ &=&  2 $As(OH)_{3,aq}$ + $CO_{2,aq}$ + $H_2O$ \\
4 $FeOOHs,fhy$ +  $CH_2O$ + 8 $H^+$ &=&  4 $Fe^{+2}$ +  $CO_{2,aq}$ + 7 $H_2O$ \\
4 $FeOOHs,gte$ +  $CH_2O$ + 8 $H^+$ &=&  4 $Fe^{+2}$ +  $CO_{2,aq}$ + 7 $H_2O$ \\
  $SO_4^{-2}$ +  2 $CH_2O$ + $H^+$  &=&   2 $CO_{2,aq}$ +  $HS^-$ + 2 $H_2O$ \\
2 $CH_2O$                           &=& $CH_4.aq$ + $CO_{2,aq}$ \\
\bottomrule
\label{tab:rdx4}
\end{tabular}
\end{table}	

The redox titration will be conducted on a suspension with the composition listed in Table \ref{tab:rdx3}.

% OK UP TO HERE

\begin{itemize}
\item Calculate the number of $\mu$moles of $CH_2O$ required to deplete each 
of the oxidants listed in Table \ref{tab:rdx3}.  
Record the number of $\mu$moles of $CH_2O$ required to reduce each of the oxidants 
and the running total of $CH_2O$ in Table \ref{tab:rdx6}.  
These will be used to set up the PHREEQC input file for the redox titration. 

\item Set up PHREEQC input file. Use the Excel file provided to set up the input file.  
This file has thermodynamic data for arsenic that are not in the Phreeqc.dat/Amm.dat database, 
an alternative to the $N$ called $Nit$, 
and phases $FeOOH(s)$ and $Mangox(s)$ defined so that their 
logK values can be changed without extensive editing of the input file.  
It also has a {\color{red} \textbf{SELECTED\_OUTPUT}} 
block set up to produce a file called {\color{blue} Redox\_titration.sel}.

\end{itemize}

\begin{table}[h]
\caption{Mass of $CH_2O$ required to titrate each oxidant in Table \ref{tab:rdx3}}
\begin{center}
\begin{tabular}{lccc} 
\hline

EA	        & Initial         & $\mu$moles $CH_2O$ 	 & Total $\mu$moles  \\
            &   $\mu$moles    &  to deplete EA       &  $CH_2O$ added    \\ \hline
$O_2$	      &     250	        &       250	           &     250        \\
$NO_3^-$    &	1500            &                      &                \\		
$MnO_{2,s}$	& 800	            &                      &                \\	
$As(5)$	    & 1               &                      &                \\		
$FeOOH_s$	  & 8000            &                      &                \\		
$U(6)$      &	1               &                      &                \\		
$SO_4^{-2}$	& 260             &                      &                \\	 \hline

\end{tabular}
\label{tab:rdx6}
\end{center}
\end{table}


\begin{enumerate}
\item Enter the groundwater composition from Table \ref{tab:rdx3}, 
including the aqueous oxidants under {\color{red} \textbf{SOLUTION 1}}.  
Note that $Nit(5)$ is used for nitrate in place of $N(5)$.  
This is defined as a {\color{red} \textbf{SOLUTION\_MASTER\_SPECIES}} and the thermodynamic 
data for the species are supplied under {\color{red} \textbf{SOLUTION\_SPECIES}}.  
Use the $Cl^-$ concentration to charge-balance the solution.

\item Enter the solid-phase oxidants from Table \ref{tab:rdx3}, 
including the phase-name, saturation index, and number of moles initially present 
under {\color{red} \textbf{EQUILIBRIUM\_PHASES 1}}.  
Use $FeOOHs$ as the name for hydrous ferric oxide.  
Thermodynamic data for a phase by this name is provided 
under the {\color{red} \textbf{PHASES}} input data block provided.  
In addition to the oxidants from Table \ref{tab:rdx1}, include $FeS(ppt)$ (amorphous FeS) 
as a phase the can form if it becomes supersaturated.

\item The {\color{red} \textbf{SURFACE 1}} data block provides 
the data needed to include surface complexation sites on $FeOOH(s)$.  
The sites are linked to the {\color{red} \textbf{EQUILIBRIUM\_PHASE}} $FeOOH(s)$ 
so that as it dissolves, surface complexation sites will disappear.  
{\color{red} \textbf{equilibrate 1}} forces the initial surface assemblage 
(species bound to the surface sites) to be in equilibrium with {\color{red} \textbf{Solution 1}}.

\item The titration will be conducted using the {\color{red} \textbf{REACTION n}} keyword.  
As $CH_2O$ is added, $C$ will be converted to $C(4)$ or $C(-4)$, 
depending on which is thermodynamically favorable under the conditions in the simulations.  
The total number of $\mu$moles of $CH_2O$ will be added in a series of steps.  
For example, for titrating $O_2$ we know that 250 $\mu$moles 
of $CH_2O$ must be added (Table \ref{tab:rdx6}).  
Here we chose to add them in 15 steps.  
By using {\color{red} \textbf{INCREMENTAL\_REACTIONS True}}, 
results are reported at the end of each step to the output files 
(both the default and Selected\_Output files)
 
\item Next, the {\color{red} \textbf{SOLUTION}}, 
{\color{red} \textbf{EQUILIBRIUM\_PHASE}} assemblage, and {\color{red} \textbf{SURFACE}} 
assemblage after the last incremental addition of $CH_2O$ are saved using 
the {\color{red} \textbf{SAVE}} key word.  
Test the file to make sure it runs without error(s).  
Open the selected output file to make sure the O2 has been titrated away.

\item Next, set up the part of the redox titration where nitrate is completely 
removed by incremental additions of $CH_2O$.  
You have calculated the total number of $\mu$moles of $CH_2O$ 
to add to remove the nitrate in Table 6.5. 
Using the {\color{red} \textbf{USE}} keyword, 
retrieve the compositions of the {\color{red} \textbf{SOLUTION}}, 
{\color{red} \textbf{EQUILIBRIUM\_PHASE}} assemblage, and {\color{red} \textbf{SURFACE}} 
assemblages that were saved at the end of the $O_2$ titration.  
Use the {\color{red} \textbf{REACTION}} keyword (with {\color{red} \textbf{INCREMENTAL\_STEPS}}) 
to add the number of $\mu$moles of $CH_2O$ required to reduce all of the 
nitrate in a reasonable number of steps (for example, 20).  
Don’t forget to save the compositions of the 
{\color{red} \textbf{SOLUTION}}, {\color{red} \textbf{EQUILIBRIUM\_PHASE}} 
assemblage, and {\color{red} \textbf{SURFACE}} assemblages for use in the subsequent simulations.

\item Follow the same steps to set up titrations for $Mangox(s)$, $FeOOH(s)$, and $S(6)$.  
You could try with 20 steps to titrate nitrate ($Nit(5)$), 
10 steps to titrate $Mangox(s)$, 20 steps to titrate $FeOOH(s)$, 
and 35 steps to titrate sulfate.

\end{enumerate}

\begin{itemize}

\item Peruse the selected output file for $O(0)$, $Nit(5)$, etc.  
How does the sequence compare with what you expected?  
The results are plotted in 
Figures \ref{fig:doug_redox_1} - \ref{fig:doug_redox_4}.
 
\item This simulation assumed that the logK for $FeOOH(s)$ was 4.89, 
which applies to $ferrihydrite$. What happens if this is changed to the logK for $goethite$, i.e., -1.0 ?

\end{itemize}

\begin{figure}[h]
\centering
\includegraphics[width=1.0\textwidth, angle = 0]{doug_redox_1}
\caption[]{}
\label{fig:doug_redox_1}
\end{figure}

\begin{figure}[h]
\centering
\includegraphics[width=1.0\textwidth, angle = 0]{doug_redox_2}
\caption[]{}
\label{fig:doug_redox_2}
\end{figure}

\begin{figure}[h]
\centering
\includegraphics[width=1.0\textwidth, angle = 0]{doug_redox_3}
\caption[]{}
\label{fig:doug_redox_3}
\end{figure}

\begin{figure}[h]
\centering
\includegraphics[width=1.0\textwidth, angle = 0]{doug_redox_4}
\caption[]{}
\label{fig:doug_redox_4}
\end{figure}